\scanpage{166}
\mspage{165}
VII.17

elements in a minimal basis of $\uline M$ is equal to the dimension of $R$.

A noetherian domain $R$ is said to be a \uline{regular domain} if and
only if $R_{\mathfrak p}$ is a regular local domain for every prime ideal $\mathfrak p$ of $R$.
(A necessary and sufficient condition that $R$ be a regular domain is
that $R_{\uline M}$ be a regular local domain for every maximal ideal $\uline M$). If
$R$ is a regular domain $R=\bigcap_{\{\uline M\}}R_{\uline M}$ where $\{\uline M\}$ is the set of
maximal ideals of $R$.

A \uline{semi-local ring} is a noetherian ring with a finite number of
maximal ideals.

We will sketch the proofs of the following lemmas, which are
well-known facts in the theory of local rings.

LEMMA VII.14. If $R$ is a regular local domain and $S$ is an
unramified extension, then $S$ is a semi-local regular domain.

\uline{Proof:} Since $S$ is unramified over $R$, $S$ is regular, by a
special case of [2, Corollary 5.2]. Furthermore, any finite
extension of a local domain is semi-local. We see this by considering
$R/\uline M\otimes_R S=S/\uline M S$ where $\uline M$ is the maximal ideal of $R$. $S/\uline M S$ is
finitely generated over the field $R/\uline M$ and therefore has but a finite
number of maximal ideals. But the maximal ideals of $S$ are in one-
to-one correspondence with the maximal ideals of $S/\uline M S$. Thus $S$
is semi-local.
